% !TEX program = xelatex
% Báo cáo survey bài báo ILCHO
% Biên dịch bằng XeLaTeX (cần font hỗ trợ tiếng Việt, vd. Times New Roman có sẵn trên Windows).
\documentclass[12pt,a4paper]{article}

\usepackage{fontspec}
\setmainfont{Times New Roman}

\usepackage[a4paper,margin=2.5cm]{geometry}
\usepackage{booktabs}
\usepackage{array}
\usepackage{enumitem}
\usepackage{caption}
\renewcommand{\tablename}{Bảng}
\usepackage{titlesec}
\usepackage{setspace}
\onehalfspacing

\titleformat{\section}{\large\bfseries}{\thesection.}{0.5em}{}
\titleformat{\subsection}{\normalsize\bfseries}{\thesubsection.}{0.5em}{}

\setlength{\parindent}{0pt}
\setlength{\parskip}{6pt}

\title{\textbf{BÁO CÁO SURVEY BÀI BÁO}\\[4pt]
\large ``Intelligent Handover Scheme for Improved 6G NTN\\
LEO Satellite Network Performance'' (ILCHO)}
\author{Sinh viên: ........................................\\
Giảng viên hướng dẫn: Cô ........................................}
\date{}

\begin{document}
\maketitle
\vspace{-2em}

\section{Thông tin bài báo}

M. Choi, M. Park, J. Kim, J.-M. Chung, \emph{``Intelligent Handover Scheme for Improved 6G
NTN LEO Satellite Network Performance,''} IEEE Transactions on Mobile Computing, vol. 25,
no. 5, 05/2026, tr. 6863--6880. DOI: 10.1109/TMC.2025.3642278.

\section{Bối cảnh \& vấn đề đặt ra}

6G hướng tới phủ sóng toàn cầu, trong đó mạng phi mặt đất (NTN) dùng vệ tinh quỹ đạo thấp
(LEO, cao độ 300--2.000 km) là mảnh ghép quan trọng để phủ vùng thưa dân và tăng dung lượng.
Do bay nhanh (hàng km/s) và vùng phủ hẹp, một thiết bị người dùng (UE) trong chòm vệ tinh
LEO mật độ cao (kiểu Starlink) phải \textbf{chuyển giao (handover) liên tục} — theo 3GPP TR
38.821, mỗi 6,6--132,3 giây tuỳ vị trí. Đo cường độ tín hiệu (RSRP) kiểu mạng mặt đất không
đáng tin cậy ở đây vì các vệ tinh ở xa và tín hiệu giữa chúng gần giống nhau. 3GPP đề xuất cơ
chế \textbf{Conditional Handover (CHO, Rel-16)} — chuẩn bị trước nhiều ứng viên, chuyển giao
chủ động khi điều kiện thoả thay vì chờ đo đạc — nhưng \textbf{chưa giải quyết} bài toán chọn
vệ tinh đích nào là tốt nhất khi có hàng chục/hàng trăm UE cùng cạnh tranh tài nguyên.

\section{Đóng góp chính của bài báo}

Bài báo đề xuất \textbf{ILCHO} = CHO + học tăng cường đa tác tử (\textbf{QMIX}) để chọn vệ
tinh đích, với 4 đóng góp:

\begin{enumerate}[leftmargin=1.3em]
  \item \textbf{Nền tảng giải tích (3 Lemma):} chứng minh vị trí vệ tinh là hàm \emph{tất
        định} theo thời gian (từ tham số quỹ đạo Walker-Delta) $\Rightarrow$ suy ra được điều
        kiện ``khả kiến'' (accessible) của mỗi UE với từng vệ tinh mà không cần đo đạc thời
        gian thực; và chứng minh bài toán gán tối ưu là \textbf{NP-hard} (quy về generalized
        assignment problem) $\Rightarrow$ biện minh cho việc dùng học máy thay vì giải chính
        xác.
  \item \textbf{Môi trường thực tế:} mô phỏng trên tham số Starlink thật (hồ sơ FCC) — Phase
        1-a ($\sim$550 km), Phase 2-a ($\sim$340 km), và kịch bản hybrid — thay vì mô hình giả
        định dưới 100 vệ tinh như phần lớn công trình trước.
  \item \textbf{Công thức MARL khả mở rộng:} không gian hành động của mỗi agent chỉ là số vệ
        tinh \emph{khả kiến tối đa} $N_{\max}$ (suy từ Lemma 2), không phải tổng số vệ tinh
        toàn hệ thống $\Rightarrow$ độ phức tạp huấn luyện/suy luận \textbf{tuyến tính theo số
        UE}, không phụ thuộc tổng số vệ tinh — đây là điểm kỹ thuật giá trị nhất của bài báo.
  \item \textbf{Kiểm chứng thực nghiệm:} theo nhóm tác giả, đây là nghiên cứu đầu tiên kiểm
        chứng CHO điều khiển bằng RL trên một mega-constellation LEO thương mại thực thụ.
\end{enumerate}

\section{Công trình liên quan \& 4 phương pháp so sánh (baseline)}

Bài báo xếp các công trình liên quan theo hướng tiếp cận (Bảng~\ref{tab:related}) và chọn 4
phương pháp so sánh, đại diện cho 2 nhóm — heuristic đơn tiêu chí và MARL không cắt không
gian hành động — để chỉ ra cả hai đều ``sụp'' khi số UE tăng.

\begin{table}[h]
\centering\small
\caption{Các hướng tiếp cận liên quan và hạn chế được bài báo chỉ ra}
\label{tab:related}
\begin{tabular}{@{}p{3.1cm}p{5cm}p{6cm}@{}}
\toprule
\textbf{Hướng} & \textbf{Ý tưởng} & \textbf{Hạn chế (theo bài báo)} \\
\midrule
Lý thuyết trò chơi & Tối đa lợi ích UE bằng potential game & Không học, giả định lý tưởng \\
Network-flow (\textbf{HSNF}) & Đồ thị luồng, chi phí = throughput + băng thông còn lại & Tính chi phí mọi cặp UE--vệ tinh thời gian thực bất khả thi ở quy mô mega \\
RL đa tác tử độc lập (\textbf{LBSH}) & Q-learning độc lập, thưởng = RVT + tải + phạt HO & Không gian hành động nổ khi ghép vào Starlink, khó hội tụ \\
DQN + OneWeb & Tối đa throughput & Bỏ qua chi phí báo hiệu do HO thừa \\
DRL vệ tinh tự quyết & Vệ tinh nguồn tự chọn, UE không gửi báo cáo đo & Chỉ 3 mặt phẳng quỹ đạo $\Rightarrow$ ``curse of dimensionality'' khi lên chòm thật \\
\bottomrule
\end{tabular}
\end{table}

\textbf{4 baseline dùng để so sánh với ILCHO:} MD-CHO (chọn vệ tinh khả kiến gần nhất),
MVT-CHO (chọn vệ tinh có thời gian nhìn thấy còn lại — RVT — dài nhất), HSNF (gán tham lam
tối đa hoá thông lượng tức thời có ràng buộc dung lượng), LBSH (RL đa tác tử độc lập + thưởng
cân bằng tải). \textbf{Thông điệp so sánh:} heuristic đơn tiêu chí (MD/MVT) sụp vì nhiều UE
gần nhau cùng chọn một vệ tinh ``tốt nhất'' theo tiêu chí $\Rightarrow$ nghẽn; LBSH sụp vì
không gian hành động không được cắt gọn; ILCHO ổn định nhờ cả ba: thưởng đa mục tiêu, không
gian hành động $=N_{\max}$, và QMIX xử lý tương tác nhiều agent.

\section{Phương pháp}

\subsection{Mô hình hoá bài toán}

Mỗi UE là một agent; quan sát cục bộ của agent $k$ tại bước thời gian $t$ là danh sách $N_{\max}$
vệ tinh khả kiến, mỗi vệ tinh mang 4 đặc trưng: chỉ số vệ tinh, khoảng cách tới nó, số kênh nó
đang dùng, và RVT (thời gian còn nhìn thấy được). Hành động là chọn 1 trong $N_{\max}$ vệ tinh
đó làm vệ tinh đích (T-LEO). Đây là bài toán \textbf{ra quyết định tuần tự có nhiều tác tử},
không phải bài toán dự đoán: vị trí vệ tinh được suy ra tất định từ tham số quỹ đạo (Lemma 1),
không phải học từ dữ liệu lịch sử.

\subsection{Hàm thưởng — điểm nhấn kỹ thuật của bài báo}

Thưởng phạt nặng khi chuyển giao thất bại, thưởng nhẹ khi chuyển giao thành công; ngoài ra là
một hàm chi phí \textbf{dạng sigmoid} kết hợp thông lượng và RVT: thông lượng cao $\to$ thưởng
cao, RVT quá thấp $\to$ bị phạt. Bài báo lập luận thông lượng theo thời gian của một vệ tinh
\emph{không đơn điệu} (tăng khi vệ tinh tiến gần thiên đỉnh rồi giảm khi ra xa) trong khi RVT
giảm đơn điệu $\Rightarrow$ một hàm thưởng \emph{tuyến tính} đơn giản dễ định giá sai các vệ
tinh ở hai đầu; hàm sigmoid tạo ra một ``điểm ngọt'' tránh được vấn đề này. Bài báo có hẳn một
thí nghiệm đối chứng (huấn luyện với thưởng tuyến tính) để minh hoạ lập luận này.

\subsection{Bộ điều khiển QMIX}

Bài báo chọn QMIX (so với MADDPG, MAPPO) vì học hiệu quả và hội tụ nhanh hơn trong bài toán
hợp tác này. Mỗi agent có một mạng nơ-ron hồi quy (DRQN, dùng GRU để nhớ lịch sử quan sát) ước
lượng giá trị hành động cục bộ; một \emph{mixing network} trộn các giá trị cục bộ thành giá
trị toàn cục theo ràng buộc đơn điệu (hành động tham lam cục bộ luôn tăng giá trị toàn cục) —
đây là cách QMIX khắc phục vấn đề \emph{non-stationarity} (môi trường ``không đứng yên'' dưới
góc nhìn một agent, vì các agent khác cũng đang học) mà các phương pháp học độc lập như IQL
(dùng trong LBSH) gặp phải. Huấn luyện tập trung, thực thi phân tán (CTDE): lúc suy luận, mỗi
UE/vệ tinh chỉ cần quan sát cục bộ.

\subsection{Quy trình chuyển giao 3 pha}

ILCHO chỉ can thiệp vào \textbf{pha chuẩn bị} của CHO (dùng QMIX chọn vệ tinh đích); \textbf{pha
thực thi} vẫn dùng một sự kiện cố định dựa trên khoảng cách (chuyển giao khi vệ tinh đích gần
hơn vệ tinh đang phục vụ một khoảng bù \emph{offset}); \textbf{pha hoàn tất} giống chuyển giao
truyền thống.

\section{Kết quả thực nghiệm bài báo công bố}

Thiết lập: Starlink Phase 1-a/2-a/hybrid, vùng dịch vụ hẹp [39--41\textdegree{}N], 5--100 UE,
huấn luyện 3.000--30.000 episode. Các kết quả chính (Bảng~\ref{tab:results}):

\begin{table}[h]
\centering\small
\caption{Tóm tắt kết quả chính bài báo công bố (Starlink Phase 2-a, 5$\to$100 UE)}
\label{tab:results}
\begin{tabular}{@{}p{3.6cm}p{3.6cm}p{3cm}p{3cm}@{}}
\toprule
\textbf{Chỉ số} & \textbf{ILCHO} & \textbf{Baseline tốt nhất} & \textbf{Baseline kém nhất} \\
\midrule
Số chuyển giao / UE & 10,8 $\to$ 13,0 (ổn định) & $\sim$10 (HSNF, ổn định) & 10 $\to$ 38 (MD-CHO) \\
Số chuyển giao thất bại / UE & 0,48 $\to$ 1,38 & -- & tăng mạnh, mọi baseline \\
Thông lượng (bps/Hz) & \textbf{3,85 $\to$ 3,76} (cao nhất) & $\sim$3,65--3,70 (HSNF) & $\sim$3,4--3,6 \\
\bottomrule
\end{tabular}
\end{table}

Ở 100 UE, ILCHO giữ số chuyển giao gần như không đổi (mỗi $\sim$46 giây/lần) dù tải tăng 20
lần so với 5 UE, trong khi mọi baseline tăng vọt. Độ công bằng (Jain's Fairness Index) đạt
\textbf{0,9798}; phương sai chiếm dụng kênh (đo cân bằng tải) của ILCHO (0,072) thấp hơn hẳn
MD-CHO (0,372) và MVT-CHO (0,400) — dù hàm thưởng không có số hạng cân bằng tải tường minh,
đây là hệ quả gián tiếp của việc phạt nặng chuyển giao thất bại (agent tự học né vệ tinh quá
tải).

\section{Điểm mạnh}

\begin{itemize}[leftmargin=1.2em]
  \item Cơ sở giải tích chặt chẽ, mô hình dựng từ tham số một chòm vệ tinh thương mại thật.
  \item Thiết kế không gian hành động theo $N_{\max}$ giải quyết đúng vấn đề mà các công trình
        MARL trước đó vấp phải khi mở rộng lên quy mô mega-constellation.
  \item So sánh công bằng với nhiều baseline, nhiều kịch bản chòm vệ tinh.
  \item Có thí nghiệm đối chứng (ablation) cho lựa chọn thiết kế hàm thưởng, có lập luận vật lý
        rõ ràng đi kèm.
\end{itemize}

\section{Hạn chế}

\begin{itemize}[leftmargin=1.2em]
  \item \textbf{Không phải ``dự đoán'' theo nghĩa học máy.} Vị trí vệ tinh là tất định (Lemma
        1); phần ``trí tuệ'' nằm ở \emph{ra quyết định}, không phải học một đại lượng bất
        định (RSRP tương lai, tải tương lai...). Câu chữ ``movement predictions'' ở phần kết
        luận của bài báo dễ gây hiểu nhầm.
  \item UE được giả định \textbf{cố định} (VSAT) — không có UE di động (điện thoại, tàu, máy
        bay).
  \item \textbf{Bỏ qua hiệu ứng Doppler} — chỉ giả định đã được bù ở lớp vật lý.
  \item Huấn luyện tốn kém (3.000--30.000 episode), chưa phân tích khả thi học trực tuyến
        (online learning) trên bộ điều khiển thời gian thực.
  \item Vùng địa lý thử nghiệm hẹp (một dải vĩ độ), chưa quét theo vĩ độ cao/vùng cực.
  \item \textbf{Thiếu chi tiết để tái lập}: không công bố tần số sóng mang chính xác, giá trị
        offset khoảng cách, các hằng số của hàm thưởng, hệ số pha chòm vệ tinh, siêu tham số
        gốc của 2 baseline HSNF/LBSH, và không báo cáo độ lệch chuẩn/số lần chạy lặp lại.
\end{itemize}

\section{Khoảng trống nghiên cứu \& hướng phát triển có thể khai thác}

Đặt bài báo vào bức tranh chung ``AI cho handover'', có thể phân loại theo trục
\textbf{dự đoán} (prediction — học một đại lượng bất định từ dữ liệu) và \textbf{ra quyết
định} (decision/control — như ILCHO). ILCHO thuần thuộc nhóm ra quyết định. Một số hướng mở:

\begin{itemize}[leftmargin=1.2em]
  \item \textbf{Thêm thành phần dự đoán thật sự} (LSTM/Transformer dự báo SINR hiệu dụng
        tương lai, hoặc xác suất chuyển giao thất bại) làm điều kiện chuẩn bị/thực thi CHO,
        khác biệt rõ với cách tiếp cận thuần hình học của ILCHO.
  \item Mở rộng sang \textbf{UE di động} và \textbf{bù Doppler không hoàn hảo}.
  \item \textbf{Học trực tuyến / liên tục} thay vì huấn luyện offline một lần.
  \item \textbf{Multi-connectivity} (chuẩn bị nhiều kết nối thay vì hard handover).
  \item Xây \textbf{benchmark chuẩn} cho bài toán handover trong NTN (hiện mỗi công trình dùng
        một bộ mô phỏng, tham số riêng, khó so sánh trực tiếp).
\end{itemize}

\section{Kết luận}

ILCHO là một cách tiếp cận có cơ sở kỹ thuật vững cho bài toán chọn vệ tinh đích trong chuyển
giao NTN, với đóng góp cốt lõi là thiết kế không gian hành động giúp học tăng cường đa tác tử
mở rộng được tới quy mô mega-constellation — điều các công trình MARL trước đó chưa làm được.
Hạn chế lớn nhất về mặt học thuật là cách dùng từ ``dự đoán'' không chính xác (bài toán thực
chất là ra quyết định trên thông tin tất định), và về mặt thực nghiệm là thiếu nhiều chi tiết
cần thiết để tái lập đầy đủ kết quả.

\end{document}
